% related work: drafted 2026-10-09 from the notes in literature/ (one notes file per cited work).
% Not yet \input in paper.tex; the coordinator places it.
\section{Related work}\label{sec:related}

\emph{Memetics and cultural evolution.} Dawkins proposed that ideas, tunes and practices are replicators of their own, which he called memes, and that memes which do well together form co-adapted complexes~\cite{dawkins1976}. In this view a meme is selected for being copied, not for helping its carriers. Boyd and Richerson built population models of cultural change, in which variants spread through content, conformist and prestige biases in social learning~\cite{boyd1985}. Sperber's school argues that transmission is mostly reconstruction, not copying: learners rebuild what they receive, and shared biases pull variants toward attractors~\cite{claidiere2014}. These frames describe how the frequency of a variant changes. None of them asks whether a set of variants is one thing, with dynamics that its parts do not have. A memeplex in memetics is a list of traits that occur together. No test separates it from a list that occurs together by chance or from a common cause. Condition~(1) of Sec.~\ref{sec:definition} is such a test. Cultural attraction also supplies the main alternative to our hypothesis. A practice that every model falls into without reading anyone is an attractor of the models' own biases, not an ideology carried by hosts.

\emph{Group agents, superorganisms and institutions.} List and Pettit argue that some groups are agents~\cite{list2011}. The attitudes of such a group are fixed by what its members do, but they need not match any member's attitudes or the majority's. The group keeps its identity while its members change, because its procedures persist. H\"olldobler and Wilson describe insect colonies as superorganisms, with division of labor, communication and colony-level homeostasis inside a bounded membership~\cite{holldobler2009}. North defines institutions as the formal and informal rules that structure interaction; they persist and are enforced while the people who follow them turn over~\cite{north1991}. These works give our two reference cases. The colony is a fixed group with a group state. The government is a set of rules that runs over changing people. The same works give marks to look for: division of labor, repair, and the enforcement of norms. They do not give a statistic. Schwitzgebel adds the point we rely on most: nothing forbids a group entity from being made of parts that are themselves agents~\cite{schwitzgebel2015}.

\emph{Measurable criteria for individuals and agents.} Bertschinger \emph{et al.} measured the autonomy of a system as the information its present state carries about its next state, given the history of its environment~\cite{bertschinger2008}. Krakauer \emph{et al.} turned this into a theory of individuality~\cite{krakauer2020}. A candidate system is an individual to the degree that it carries information from its own past into its own future beyond its environment. Its boundary is found by growing the system while this information rises. Kolchinsky and Wolpert add a test of agency: information is semantic if scrambling it lowers the system's viability~\cite{kolchinsky2018}. These quantities are defined for any split into system and environment, including a pattern and the agents that carry it. That is why we use them. Causal-emergence measures ask a related question of macroscopic variables: does the macro state predict the future beyond every part~\cite{rosas2020,mediano2022}? Biehl \emph{et al.} point out that a fixed set of variables cannot follow an entity whose material changes~\cite{biehl2016}. They define entities as spatiotemporal patterns of values instead. An ideology on changing hosts is such an entity. We handle it by taking the state of the pattern, not the state of its hosts, as the system variable. Kirchhoff \emph{et al.} describe living systems as Markov blankets nested inside larger blankets~\cite{kirchhoff2018}. Levin sizes an agent by the spatial and temporal reach of the goals it pursues~\cite{levin2022}. Both allow agents inside agents. Integrated information theory takes the opposite view: its exclusion postulate keeps only the most integrated of overlapping systems~\cite{oizumi2014}. We do not impose exclusion. We report every local maximum of individuality, at the level of the agent and above it.

\emph{Culture in populations of language-model agents.} Generative agents in a sandbox town spread a planted piece of news to half the population in two simulated days~\cite{park2023}. Populations of language models in a naming game reach shared conventions. They show collective biases that no single agent has, and committed minorities can flip them~\cite{ashery2025}. Chains of language models that pass text along drift toward attractors that depend on the model and on how open the task is~\cite{perez2025}. Societies of language-model agents in a donor game evolve norms of indirect reciprocity across generations, or fail to, depending on the base model~\cite{vallinder2024}. These studies are experiments. Each designs its setting: one model per population, a fixed game or task, explicit generations, and an event or payoff that the authors choose. This control lets them measure transmission cleanly. It also limits the norms they can find to the norms the design admits. Our system is one open deployment. It holds 21 to 32 frontier models from several labs, with private roles, real tools, months of runtime, and erasure of working memory every few dozen calls. Nobody planted an event or defined a payoff. What we lose in control we must recover with nulls: the operator's fields and the role texts in the environment $E$, frequency-matched pseudo-patterns, read versus unread exposure, and natural experiments such as memory wipes and the arrival of newcomers.

\emph{Position of this paper.} The paper joins two lines of work. From cultural evolution it takes the object: a pattern of ideas and practices on a population of carriers. From the information theory of individuality and agency it takes the test. We know of no earlier work that measures individuality or semantic information at the level of a cultural pattern, in people or in machines.
